\chapter{TCP: la gestione della connessione}

% ══════════════════════════════════════════════════════════════
\section{Il problema}
% ══════════════════════════════════════════════════════════════

Poiché TCP è orientato alla connessione, i due host devono mettersi d'accordo sia quando la connessione
si \textbf{apre} sia quando si \textbf{chiude}. Ma i messaggi possono perdersi o danneggiarsi, e trovare un
accordo è tutt'altro che banale: se un messaggio si perde, un lato può ritenere la connessione aperta (o
chiusa) mentre l'altro no.

Per evitare, o meglio per \textbf{mitigare}, il problema TCP usa una procedura detta \textbf{three-way
handshake} (stretta di mano a tre vie): le richieste di apertura e di chiusura devono essere
\textbf{confermate} prima che la connessione sia considerata stabilita o rilasciata. Niente conta finché non è stato
confermato, e l'apertura è costruita su questa regola.

Delle cose che fa il trasporto, questa è quella che il nostro codice ha chiesto fin dall'inizio: \texttt{connect},
\texttt{accept}, \texttt{close}. E non serve niente di nuovo nell'header: bastano tre bit di flag mai usati finora.

\begin{itemize}
    \item \textbf{SYN}: questo segmento apre una connessione. È acceso nei primi due segmenti di una connessione e mai
          più.
    \item \textbf{FIN}: il mittente non ha più dati. Chiude \textbf{una direzione}, e si invia una volta per direzione.
    \item \textbf{RST}: qui non c'è nessuna connessione. Butta via il segmento e non rimandarlo.
\end{itemize}
I campi di sequenza e di acknowledgment portano l'accordo; questi tre bit dicono di che tipo di accordo si tratta.

% ══════════════════════════════════════════════════════════════
\section{Il problema dei due eserciti}
% ══════════════════════════════════════════════════════════════

\begin{figure}[H]
\centering
\begin{tikzpicture}[every node/.style={font=\footnotesize\sffamily}]
  \fill[green!15] (-1,-1.2) -- (1.2,1.6) -- (3,0.4) -- (5,1.7) -- (7.4,-1.2) -- cycle;
  \fill[green!30] (-1,-1.2) -- (0.6,0.4) -- (2,-0.3) -- (3.2,-0.9) -- (4.8,-0.2) -- (6.2,0.5) -- (7.4,-1.2) -- cycle;
  \node[draw, fill=blue!40, rounded corners, minimum width=1.5cm] (b1) at (0.4,0.9) {blu \#1};
  \node[draw, fill=blue!40, rounded corners, minimum width=1.5cm] (b2) at (6,1.0) {blu \#2};
  \node[draw, fill=white, rounded corners, minimum width=2.6cm, minimum height=0.8cm] (w) at (3.2,-0.6) {esercito bianco};
  \draw[msg, dashed, netred, thick] (b1) to[bend right=15] node[below, sloped, text=netred] {messaggero} (b2);
  \node[text width=6.5cm, anchor=west] at (8,0.2) {Il bianco è più forte di ciascun esercito blu, ma più debole dei due insieme: i blu vincono solo se attaccano \textbf{nello stesso momento}. Per accordarsi possono mandare messaggeri attraverso la valle, dove possono essere catturati.};
\end{tikzpicture}
\caption{Il problema dei due eserciti: comunicare attraverso un canale inaffidabile.}
\end{figure}

Un esercito bianco è accampato in una valle; sulle due colline ai lati ci sono due eserciti blu. Il
bianco è più grande di ciascuno dei due eserciti blu, ma insieme i blu sono più numerosi: vincono
\textbf{solo se attaccano contemporaneamente}. Per sincronizzarsi devono mandare messaggeri attraverso la
valle, dove possono essere catturati (comunicazione \textbf{inaffidabile}). Esiste un protocollo che permetta
ai blu di vincere?

\begin{figure}[H]
\centering
\begin{tikzpicture}[yscale=0.8, every node/.style={font=\scriptsize\sffamily}]
  \begin{scope}
    \node[font=\footnotesize\bfseries\sffamily] at (1.9,1.1) {Due vie};
    \timeline{a}{0}{blu \#1}{3.4}
    \timeline{b}{3.8}{blu \#2}{3.4}
    \draw[msg] (0,-0.3) -- node[above, sloped] {``attacchiamo oggi, ok?''} (3.8,-1.1);
    \draw[msg, netgreen] (3.8,-1.4) -- node[above, sloped] {``d'accordo!''} (0,-2.2);
    \node[text=netred, anchor=west] at (3.9,-2.3) {sarà arrivato?};
  \end{scope}
  \begin{scope}[xshift=7.5cm]
    \node[font=\footnotesize\bfseries\sffamily] at (1.9,1.1) {Tre vie};
    \timeline{a}{0}{blu \#1}{3.4}
    \timeline{b}{3.8}{blu \#2}{3.4}
    \draw[msg] (0,-0.3) -- node[above, sloped] {``attacchiamo oggi, ok?''} (3.8,-1.1);
    \draw[msg, netgreen] (3.8,-1.4) -- node[above, sloped] {``d'accordo!''} (0,-2.2);
    \draw[msg] (0,-2.4) -- node[above, sloped] {``d'accordo!''} (3.8,-3.2);
    \node[text=netred, anchor=east] at (-0.1,-2.9) {e il mio, sarà arrivato?};
  \end{scope}
\end{tikzpicture}
\caption{Aggiungere conferme sposta solo il dubbio sull'ultimo messaggio.}
\end{figure}

\begin{itemize}
    \item Il comandante del blu \#1 manda: ``\emph{Propongo di attaccare oggi, va bene?}''. Il messaggio arriva, il
          blu \#2 è d'accordo e la sua risposta torna sana e salva. L'attacco ci sarà? Probabilmente no: il
          comandante \#2 \textbf{non sa} se la sua risposta è arrivata, e se non fosse arrivata il \#1 non attaccherebbe.
    \item Proviamo con tre vie: il \#1 conferma la risposta. Se nessun messaggio si perde, il \#2 riceve la conferma,
          ma ora è il \#1 a esitare: non sa se la sua conferma è arrivata.
    \item Con quattro vie non cambia nulla. Si può \textbf{dimostrare} che \textbf{non esiste alcun protocollo} che
          risolva il problema con certezza.
\end{itemize}

\info{Imperfetto ma adeguato}{
    Il three-way handshake \textbf{non è perfetto}, ma nella pratica è di solito \textbf{più che adeguato}: le
    situazioni residue vengono gestite con timer e ritrasmissioni. L'incertezza non riguarda il messaggio: riguarda che
    cosa sa l'altro.
}

\dfn{Che cosa è impossibile}{
    Che entrambi gli estremi siano \textbf{certi, nello stesso istante}, che l'altro sia pronto. Nessun numero di messaggi
    ci riesce, su un canale che perde messaggi.
}

\begin{itemize}
    \item Ciò di cui ha bisogno TCP è più debole: ciascun estremo ha imparato il numero iniziale dell'altro, e sa che il
          proprio è stato ricevuto.
    \item Tre segmenti danno esattamente questo, e nessuno va creduto sulla parola: ciascuno è confermato dal successivo.
    \item Il dubbio che resta è quello del mittente sul suo ultimo messaggio, e lo risolve il primo segmento di dati.
    \item Dove quel dubbio costa qualcosa è alla chiusura, e ci torneremo con i timer.
\end{itemize}

% ══════════════════════════════════════════════════════════════
\section{L'apertura della connessione}
% ══════════════════════════════════════════════════════════════

\begin{enumerate}
    \item La \texttt{connect} del client invia una richiesta di connessione: un segmento \textbf{SYN}, senza dati.
    \item Il server risponde con un SYN proprio che conferma anche quello del client: un solo segmento, \textbf{SYN-ACK}
          (o SYNACK).
    \item Il client invia la conferma finale: un segmento \textbf{ACK}, che può già portare dati dell'applicazione.
\end{enumerate}
Un'intera andata e ritorno, più il tempo che il server impiega a rispondere, prima del primo byte della richiesta.

\subsection{I dettagli}

\begin{figure}[H]
\centering
\begin{tikzpicture}[yscale=0.9, every node/.style={font=\footnotesize\sffamily}]
  \timeline{c}{0}{Client}{5.2}
  \timeline{s}{8}{Server}{5.2}
  \node[anchor=east, text=black!60] at (-0.2,0) {\texttt{connect()}};
  \node[anchor=west, text=black!60] at (8.2,0) {\texttt{listen()}, \texttt{accept()}};
  \draw[msg, very thick, netblue] (0,-0.5) -- node[above, sloped, align=center] {SYN = 1, Seq = client\_isn\\(nessun dato)} (8,-1.7);
  \node[anchor=west, text width=3.4cm, align=left] at (8.2,-1.9) {alloca buffer e variabili};
  \draw[msg, very thick, netgreen] (8,-2.2) -- node[above, sloped, align=center] {SYN = 1, ACK = 1, Seq = server\_isn\\Ack = client\_isn + 1} (0,-3.4);
  \node[anchor=east, text width=3cm, align=right] at (-0.2,-3.6) {alloca buffer e variabili};
  \draw[msg, very thick, netblue] (0,-3.9) -- node[above, sloped, align=center] {SYN = 0, ACK = 1, Seq = client\_isn + 1\\Ack = server\_isn + 1 (può contenere dati)} (8,-5.1);
  \node[anchor=west, text=netgreen!60!black, font=\footnotesize\bfseries\sffamily] at (8.2,-5.1) {connessione stabilita};
\end{tikzpicture}
\caption{Il three-way handshake nel dettaglio.}
\end{figure}

\begin{enumerate}
    \item Il client invia un segmento \textbf{SYN} che ha:
    \begin{itemize}
      \item \textbf{nessun dato} applicativo;
      \item il bit \textbf{SYN a 1};
      \item come numero di sequenza un \textbf{numero iniziale casuale} (\texttt{client\_isn}, \emph{initial sequence number}).
    \end{itemize}
    \item Quando il SYN arriva (si spera), il server \textbf{alloca buffer e variabili} TCP e risponde con un segmento
          \textbf{SYNACK} che ha:
    \begin{itemize}
      \item nessun dato applicativo;
      \item i bit \textbf{SYN e ACK a 1};
      \item numero di acknowledgment \texttt{client\_isn + 1};
      \item come numero di sequenza un proprio numero iniziale casuale (\texttt{server\_isn}).
    \end{itemize}
    \item Quando il SYNACK arriva (si spera), anche il client alloca buffer e variabili e invia un segmento
          \textbf{ACK} finale che ha:
    \begin{itemize}
      \item eventualmente \textbf{dati} applicativi (per esempio la richiesta HTTP);
      \item il bit \textbf{SYN a 0} (la connessione ora è stabilita) e il bit \textbf{ACK a 1};
      \item numero di acknowledgment \texttt{server\_isn + 1}.
    \end{itemize}
\end{enumerate}

\warning{Perché il SYN ``consuma'' un numero di sequenza}{
    Il SYN non porta dati, eppure l'ACK di risposta vale \texttt{client\_isn + 1}: per convenzione SYN (e FIN)
    contano come un byte, così possono essere confermati come qualsiasi altro dato.
}

\subsection{Perché il numero iniziale è estratto a caso}

\begin{itemize}
    \item Niente obbliga una connessione a partire da zero, e partire da lì sarebbe la scelta peggiore possibile.
    \item Un segmento di una connessione precedente fra gli stessi quattro numeri può essere ancora in giro nella rete. Con un
          punto di partenza nuovo cade fuori dalla finestra della nuova connessione, e viene scartato invece di essere
          accettato come dato.
    \item Un attaccante che indovina i numeri può iniettare segmenti in una connessione che non vede: renderli casuali è una
          difesa.
    \item E il $+1$ che torna indietro è una piccola prova a sé: solo chi ha ricevuto il SYN poteva scriverlo.
\end{itemize}

\subsection{Diagramma degli stati (apertura)}

\begin{figure}[H]
\centering
\begin{tikzpicture}[every node/.style={font=\scriptsize\sffamily}, st/.style={draw, rounded corners=6pt, fill=cloudbg, thick, minimum width=2.4cm, minimum height=0.75cm, font=\footnotesize\bfseries\sffamily}]
  \node[st] (closed) at (0,0) {CLOSED};
  \node[st] (listen) at (-4,-2) {LISTEN};
  \node[st] (synsent) at (4,-2) {SYN SENT};
  \node[st] (synrcvd) at (-4,-4.3) {SYN RECEIVED};
  \node[st, fill=netgreen!20] (est) at (0,-6.3) {ESTABLISHED};
  \draw[-{Stealth}, thick] (closed) -- node[above, sloped, align=center] {apertura passiva\\\texttt{listen()}} (listen);
  \draw[-{Stealth}, thick] (closed) -- node[above, sloped, align=center] {apertura attiva: \texttt{connect()}\\invia SYN} (synsent);
  \draw[-{Stealth}, thick] (listen) -- node[left, align=right] {ricevuto SYN\\invia SYN+ACK} (synrcvd);
  \draw[-{Stealth}, thick] (synsent) -- node[right, align=left] {ricevuto SYN+ACK\\invia ACK} (est);
  \draw[-{Stealth}, thick] (synrcvd) -- node[below left, align=right] {ricevuto ACK} (est);
  \node[text=netblue, font=\footnotesize\bfseries\sffamily] at (-4,-0.9) {server};
  \node[text=netred, font=\footnotesize\bfseries\sffamily] at (4,-0.9) {client};
\end{tikzpicture}
\caption{Diagramma degli stati semplificato per l'apertura della connessione. Due percorsi dentro la stessa macchina: il
         server aspetta (apertura passiva) e risponde a un SYN con un SYNACK; il client chiede (apertura attiva) e invia il
         SYN. Si incontrano in ESTABLISHED, il server un segmento dopo il client.}
\end{figure}

\subsection{Gli stessi stati, visti dal nostro codice}

\begin{itemize}
    \item \texttt{listen} è la transizione in LISTEN. Da quell'istante il kernel risponde ai SYN \textbf{da solo}, senza
          alcun processo coinvolto.
    \item \texttt{connect} invia il SYN e si blocca; ritorna quando è arrivato il SYNACK ed è partito il terzo segmento.
    \item \texttt{accept} \textbf{non partecipa affatto} all'handshake: consegna una connessione già stabilita. Un server può
          quindi essere lento a chiamare \texttt{accept} senza che alcun client se ne accorga: la connessione è aperta, e i
          dati aspettano in un buffer.
    \item L'handshake lo fanno i due kernel. Il programma ne legge il risultato, e il suo unico resoconto è un descriptor o
          un errore.
\end{itemize}

\dfn{Le due code di listen}{
    Un socket di ascolto ha una coda delle connessioni \textbf{mezze aperte} (SYN ricevuto, terzo segmento non ancora) e
    una coda di quelle \textbf{completate} che aspettano \texttt{accept}.
}

\begin{itemize}
    \item Un SYN in arrivo crea una voce nella prima coda; l'ACK finale la sposta nella seconda.
    \item \texttt{accept} toglie una voce dalla seconda coda: il \emph{backlog} passato a \texttt{listen} è la lunghezza di
          quella coda.
    \item Se si riempie, le nuove connessioni vengono scartate o rifiutate, e il client vede un ritardo senza mai saperne il
          motivo.
    \item La prima coda è più piccola, più vecchia e sguarnita: è quella a cui mira un attaccante.
\end{itemize}

\subsection{Quanto costa l'handshake}

\begin{itemize}
    \item Prima che parta il primo byte di una richiesta, un'intera andata e ritorno è già stata spesa per mettersi
          d'accordo a parlare. Su una rete locale non è niente; attraverso un oceano è un decimo di secondo, per
          connessione.
    \item È esattamente ciò che comprano le \textbf{connessioni persistenti} del livello applicativo: si paga l'handshake
          una volta, e si riusa ciò che ha costruito.
    \item Costa anche al server: buffer e variabili vengono impegnati al primo segmento, prima che gli sia stato chiesto
          niente.
\end{itemize}

\warning{Nessuno risponde, oppure qualcuno rifiuta}{
    Il client chiama \texttt{connect} e il SYN non riceve risposta. Quanto aspetta? \textbf{Decine di secondi}: TCP
    ritrasmette il SYN qualche volta con un intervallo crescente, poi abbandona il tentativo e \texttt{connect} fallisce.
    \begin{itemize}
        \item Un rifiuto è molto più veloce: se l'host è acceso ma nessuno ha fatto \texttt{bind} su quella porta, risponde
              con un segmento con \textbf{RST}.
        \item Il programma vede due errori diversi: un timeout, o l'\texttt{ECONNREFUSED} ottenuto la prima volta che si è
              avviato il client prima del server.
        \item Per UDP quella segnalazione arrivava dal livello di rete; qui la risposta è un segmento di trasporto, dal TCP
              del peer.
    \end{itemize}
}

\subsection{Il SYN flood e i SYN cookie}

E se il terzo segmento non arriva mai? Il server ha già impegnato buffer e variabili, e tiene la connessione mezza aperta
per un minuto o più prima di recuperarla.
\begin{itemize}
    \item Inviando abbastanza SYN, la coda delle connessioni mezze aperte si riempie di voci che non si completeranno mai, e
          i client legittimi vengono rifiutati: si è esaurita una \textbf{coda}, non la banda, e l'attacco richiede pochissimo
          traffico.
    \item Gli indirizzi sorgente sono falsificati, quindi i SYNACK vanno a host che non hanno chiesto niente, e
          l'attaccante non deve mai rispondere.
    \item Il \key{SYN flood} è stato fra i primi attacchi di \emph{denial of service} documentati, ed è ancora fra i più
          comuni.
\end{itemize}

\dfn{SYN cookie}{
    La difesa capovolge il problema: \textbf{rispondere al SYN senza ricordarlo}. \texttt{server\_isn} non è casuale ma
    calcolato: un hash dei due indirizzi, delle due porte e di un segreto noto solo al server.
}

\begin{itemize}
    \item Non si tiene alcuno stato, quindi un SYN falsificato costa al server una risposta e nient'altro.
    \item Un client legittimo restituisce quel numero più uno. Il server ricalcola l'hash, confronta, e solo allora crea la
          connessione e il suo socket.
    \item La coda non si riempie mai, perché fino al terzo segmento non c'è niente dentro. Oggi lo fanno quasi tutti i
          sistemi operativi.
\end{itemize}

\begin{table}[H]
\centering
\begin{tabular}{l l l}
\toprule
\textbf{Che cosa torna} & \textbf{Che cosa significa} & \textbf{Stato della porta} \\
\midrule
SYNACK  & qualcosa ascolta lì                          & \emph{open} \\
RST     & l'host è acceso, quella porta non è servita  & \emph{closed} \\
niente  & non è mai arrivato: qualcosa l'ha scartato   & \emph{filtered} \\
\bottomrule
\end{tabular}
\caption{Le tre risposte a un SYN. Un SYN per porta, e le tre risposte mappano una macchina dall'esterno senza mai
         completare una connessione: è ciò che fa \texttt{nmap}, ed è il significato degli stati visti con le porte.}
\end{table}

% ══════════════════════════════════════════════════════════════
\section{La chiusura della connessione}
% ══════════════════════════════════════════════════════════════

La chiusura (\emph{teardown}) può essere avviata da \textbf{uno qualsiasi} dei due host. Supponiamo che a
chiudere sia il client:

\begin{figure}[H]
\centering
\begin{tikzpicture}[yscale=0.85, every node/.style={font=\footnotesize\sffamily}]
  \timeline{c}{0}{Client}{5.8}
  \timeline{s}{7}{Server}{5.8}
  \node[anchor=east, text=black!60] at (-0.2,-0.4) {\texttt{close()}};
  \draw[msg, very thick, netred] (0,-0.5) -- node[above, sloped] {FIN = 1} (7,-1.4);
  \draw[msg, very thick, netgreen] (7,-1.7) -- node[above, sloped] {ACK = 1} (0,-2.6);
  \node[anchor=west, text=black!60] at (7.2,-2.9) {\texttt{close()}};
  \draw[msg, very thick, netred] (7,-3.0) -- node[above, sloped] {FIN = 1} (0,-3.9);
  \draw[msg, very thick, netgreen] (0,-4.2) -- node[above, sloped] {ACK = 1} (7,-5.1);
  \node[anchor=west, text=netgreen!60!black, font=\footnotesize\bfseries\sffamily] at (7.2,-5.1) {chiusa};
  \draw[netorange, very thick] (-0.3,-4.2) -- (-0.3,-5.8);
  \node[anchor=east, text=netorange, text width=2.6cm, align=right] at (-0.4,-5.0) {attesa (TIME WAIT), poi chiusa};
\end{tikzpicture}
\caption{La chiusura della connessione: ACK e FIN del server possono viaggiare nello stesso segmento o in due segmenti separati.}
\end{figure}

\begin{enumerate}
    \item Il client invia un segmento speciale di chiusura: un segmento \textbf{FIN}, con il bit FIN a 1.
    \item Quando lo riceve, il server invia una conferma e, quando anche lui ha finito, un proprio segmento di
          chiusura, con ACK a 1 e FIN a 1. ACK e FIN possono viaggiare nello stesso segmento o in due separati:
          nel mezzo il server può ancora inviare dati (la connessione è chiusa solo in una direzione).
    \item Infine il client conferma il FIN del server.
\end{enumerate}

\subsection{Mezza chiusura: una direzione alla volta}

\begin{itemize}
    \item Un FIN dice «non ho più dati da inviare». Non dice «questa connessione è finita»: chiude \textbf{una direzione}.
    \item Dopo aver inviato il suo FIN il client può ancora ricevere; dopo averlo confermato il server può ancora inviare,
          per tutto il tempo che vuole.
    \item I due FIN sono quindi indipendenti, e il rilascio normale è di \textbf{quattro} segmenti: FIN, ACK, poi FIN, ACK
          nell'altro verso. Si riducono a tre solo quando il peer non ha più niente da inviare e chiude subito: allora il suo
          ACK e il suo FIN viaggiano insieme.
    \item \texttt{close} invia il FIN e rinuncia al descriptor. \texttt{shutdown(fd, SHUT\_WR)} lo invia e continua a
          leggere: è il modo in cui un client dice «questa era tutta la richiesta».
\end{itemize}

E dal lato del programma, il FIN che arriva è proprio lo \texttt{0} restituito da \texttt{recv}: un'istruzione a
chiudere, non un errore da ignorare.

\subsection{Diagramma degli stati (chiusura)}

\begin{figure}[H]
\centering
\begin{tikzpicture}[every node/.style={font=\scriptsize\sffamily}, st/.style={draw, rounded corners=6pt, fill=cloudbg, thick, minimum width=2.3cm, minimum height=0.7cm, font=\footnotesize\bfseries\sffamily}]
  \node[st, fill=netgreen!20] (est) at (0,0) {ESTABLISHED};
  \node[st] (fw1) at (-4,-1.8) {FIN WAIT 1};
  \node[st] (fw2) at (-4,-3.8) {FIN WAIT 2};
  \node[st] (closing) at (-0.6,-3.8) {CLOSING};
  \node[st] (tw) at (-4,-5.8) {TIME WAIT};
  \node[st] (cw) at (4,-1.8) {CLOSE WAIT};
  \node[st] (la) at (4,-3.8) {LAST ACK};
  \node[st, fill=lphy] (closed) at (0,-7.4) {CLOSED};
  \draw[-{Stealth}, thick] (est) -- node[above, sloped, align=center] {\texttt{close()}\\invia FIN} (fw1);
  \draw[-{Stealth}, thick] (fw1) -- node[left, align=right] {ricevuto ACK} (fw2);
  \draw[-{Stealth}, thick] (fw1) -- node[above, sloped, align=center] {ricevuto FIN\\invia ACK} (closing);
  \draw[-{Stealth}, thick] (fw2) -- node[left, align=right] {ricevuto FIN\\invia ACK} (tw);
  \draw[-{Stealth}, thick] (closing) -- node[right, align=left] {ricevuto ACK} (tw);
  \draw[-{Stealth}, thick] (tw) -- node[below left, align=right] {scade il timeout} (closed);
  \draw[-{Stealth}, thick] (est) -- node[above, sloped, align=center] {ricevuto FIN\\invia ACK} (cw);
  \draw[-{Stealth}, thick] (cw) -- node[right, align=left] {\texttt{close()}\\invia FIN} (la);
  \draw[-{Stealth}, thick] (la) -- node[below right, align=left] {ricevuto ACK} (closed);
  \node[text=netred, font=\footnotesize\bfseries\sffamily] at (-4,-0.8) {chi chiude per primo};
  \node[text=netblue, font=\footnotesize\bfseries\sffamily] at (4,-0.8) {chi chiude per secondo};
\end{tikzpicture}
\caption{Diagramma degli stati semplificato per il rilascio della connessione. Chi chiude per primo scende a sinistra;
         chi risponde scende a destra, e resta in CLOSE WAIT finché la sua applicazione non chiude. CLOSING corrisponde al
         caso raro in cui i due lati inviano il FIN contemporaneamente e nessuno dei due è secondo.}
\end{figure}

\warning{Il descriptor che nessuno ha chiuso}{
    La \texttt{recv} del server restituisce 0 e ci si dimentica di chiudere il descriptor. Che cosa si rompe? All'inizio
    \textbf{niente}. La connessione resta in CLOSE WAIT per sempre: il FIN del peer è stato confermato, ma il nostro non
    parte mai.
    \begin{itemize}
        \item Anche il peer resta in attesa. Il rilascio non finisce su nessuno dei due lati, e nessuno dei due programmi
              viene avvisato.
        \item Qualche migliaio di client dopo il server ha finito i descriptor e smette di accettare: il guasto compare
              lontano dalla sua causa.
        \item \texttt{ss -tn state close-wait} elenca esattamente quelle connessioni. Una colonna di CLOSE WAIT è
              \textbf{sempre un bug del programma}, mai della rete.
    \end{itemize}
}

\subsection{TIME WAIT}

\begin{itemize}
    \item Chi chiude per primo non passa subito a CLOSED: aspetta, di solito da trenta secondi a due minuti.
    \item \textbf{Primo motivo}: se la sua ultima conferma si perde, il peer rimanderà il suo FIN, e qualcuno deve essere lì
          a rispondere.
    \item \textbf{Secondo motivo}: segmenti di questa connessione possono essere ancora in rete, e i suoi quattro numeri
          potrebbero essere assegnati a una connessione nuova.
    \item L'attesa è circa il doppio del tempo massimo per cui un segmento può sopravvivere, così tutto ciò che appartiene
          alla vecchia connessione è morto prima.
    \item È il prezzo di chiudere per primi, e lo paga chi chiude per primo, che non è sempre il client.
\end{itemize}

\warning{Address already in use}{
    Si interrompe il server, lo si riavvia, e \texttt{bind} fallisce. Non c'è nient'altro in esecuzione. Perché?
    Perché interromperlo lo ha reso \textbf{l'estremo che ha chiuso per primo}. Le sue connessioni sono in TIME WAIT,
    ancora legate alla porta locale, e \texttt{bind} non prende una porta che una connessione viva sta usando.
    \begin{itemize}
        \item \texttt{SO\_REUSEADDR}, la riga scritta in ogni server fin dal primo, dice di fare \texttt{bind} lo stesso.
        \item È sicuro: la vecchia connessione conserva i suoi quattro numeri, e nessun suo segmento può essere consegnato al
              nuovo socket.
        \item Quando chiude prima il client, che è il caso abituale, il server non entra mai in TIME WAIT e il problema non
              compare.
    \end{itemize}
}

% ══════════════════════════════════════════════════════════════
\section{Come TCP si salva dalla perdita di pacchetti}
% ══════════════════════════════════════════════════════════════

Poiché il three-way handshake non è perfetto, TCP si affida ai \textbf{timer} per chiudere comunque le
connessioni o per ripetere le richieste. Alcuni esempi durante il rilascio:

\begin{figure}[H]
\centering
\begin{tikzpicture}[yscale=0.72, every node/.style={font=\scriptsize\sffamily}]
  \begin{scope}
    \node[font=\footnotesize\bfseries\sffamily] at (1.6,1.1) {(a) si perde l'ACK finale};
    \timeline{c}{0}{Client}{7}
    \timeline{s}{3.2}{Server}{7}
    \draw[msg] (0,-0.3) -- node[above, sloped] {FIN} (3.2,-1.0);
    \draw[msg, netgreen] (3.2,-1.2) -- node[above, sloped] {ACK + FIN} (0,-1.9);
    \draw[msglost] (0,-2.1) -- (1.7,-2.5); \node[text=netred] at (2.1,-2.2) {ACK perso};
    \draw[netorange, thick] (3.45,-1.2) -- (3.45,-3.6); \node[text=netorange, rotate=90] at (3.7,-2.4) {timeout};
    \draw[msg] (3.2,-3.6) -- node[above, sloped] {FIN (di nuovo)} (0,-4.3);
    \draw[msg, netgreen] (0,-4.5) -- node[above, sloped] {ACK} (3.2,-5.2);
    \node at (1.6,-6.2) {chiusa su entrambi i lati};
  \end{scope}
  \begin{scope}[xshift=5.6cm]
    \node[font=\footnotesize\bfseries\sffamily] at (1.6,1.1) {(b) si perde il FIN del server};
    \timeline{c}{0}{Client}{7}
    \timeline{s}{3.2}{Server}{7}
    \draw[msg] (0,-0.3) -- node[above, sloped] {FIN} (3.2,-1.0);
    \draw[msglost] (3.2,-1.2) -- (1.5,-1.6); \node[text=netred] at (1.1,-1.3) {perso};
    \draw[netorange, thick] (-0.25,-0.3) -- (-0.25,-2.6); \node[text=netorange, rotate=90] at (-0.5,-1.45) {timeout};
    \draw[msg] (0,-2.6) -- node[above, sloped] {FIN (rinviato)} (3.2,-3.3);
    \draw[msg, netgreen] (3.2,-3.5) -- node[above, sloped] {ACK + FIN} (0,-4.2);
    \draw[msg, netgreen] (0,-4.4) -- node[above, sloped] {ACK} (3.2,-5.1);
    \node at (1.6,-6.2) {chiusa su entrambi i lati};
  \end{scope}
  \begin{scope}[xshift=11.2cm]
    \node[font=\footnotesize\bfseries\sffamily] at (1.6,1.1) {(c) perdite ripetute};
    \timeline{c}{0}{Client}{7}
    \timeline{s}{3.2}{Server}{7}
    \draw[msg] (0,-0.3) -- node[above, sloped] {FIN} (3.2,-1.0);
    \draw[msglost] (3.2,-1.2) -- (1.5,-1.6);
    \draw[msg] (0,-2.4) -- node[above, sloped] {FIN} (3.2,-3.1);
    \draw[msglost] (3.2,-3.3) -- (1.5,-3.7);
    \node[text width=2.8cm, align=center] at (1.6,-5.0) {dopo $N$ timeout ciascun lato chiude comunque};
    \node[text=netred] at (0,-6.5) {chiusa}; \node[text=netred] at (3.2,-6.5) {chiusa};
  \end{scope}
\end{tikzpicture}
\caption{Timer e ritrasmissioni permettono di chiudere la connessione anche quando si perdono segmenti; dopo un certo numero di tentativi falliti la connessione viene chiusa d'ufficio.}
\end{figure}

\begin{itemize}
    \item \textbf{Si perde l'ultima conferma}: il client la invia e si considera chiuso; il server resta ad aspettarla. Il suo
          timer scade e rimanda il FIN (a cui risponde il client in TIME WAIT), oppure, finiti i tentativi, rilascia comunque
          la connessione. Non si danneggia niente: entrambi finiscono chiusi, solo non nello stesso istante, e nessuno dei
          due lo saprà mai. È il dubbio dei due eserciti, e qui costa un ritardo da una parte.
    \item \textbf{Si perde la risposta al FIN}: il client non ha saputo ancora niente. Il suo timer scade e rimanda il FIN,
          esattamente come ritrasmetterebbe un segmento di dati; il server risponde con lo stesso segmento di prima. La
          ritrasmissione alla chiusura non è un meccanismo nuovo: un timer, e una copia tenuta finché non è confermata.
    \item \textbf{Le perdite continuano}: dopo un numero limitato di tentativi il client smette e chiude da solo, e il timer
          del server chiude il suo lato. La connessione finisce senza che nessuno dei due abbia confermato niente all'altro,
          che è proprio il punto di partenza dei due eserciti. Nessun protocollo lo elimina: ciò che aggiunge TCP è un timer,
          un limite ai tentativi, e le risorse rilasciate in ogni caso.
\end{itemize}

\warning{RST: chiudere senza chiedere}{
    Un FIN termina una connessione in ordine. Un \textbf{RST} (\emph{reset}) la termina \textbf{ora}, e ciò che era ancora nei
    buffer viene scartato.
    \begin{itemize}
        \item Viene inviato ogni volta che arriva un segmento per una connessione che non esiste: nessuno in ascolto su
              quella porta, o un peer che si è riavviato. Nel caso di UDP, invece, a una porta chiusa si risponde con un
              messaggio ICMP ``port unreachable'' (lo vedremo con traceroute).
        \item Una macchina che va in crash non invia niente. Il suo peer tiene una connessione che all'altro capo non esiste
              più, e lo scopre solo quando scrive.
        \item Allora torna l'RST, e il programma lo apprende come \texttt{ECONNRESET} su una lettura, oppure
              \texttt{EPIPE} e \texttt{SIGPIPE} su una scrittura.
        \item La sonda d'inattività del kernel parte dopo circa due ore: per questo un programma che deve accorgersi di un
              peer morto invia un proprio \emph{heartbeat}.
    \end{itemize}
}

% ══════════════════════════════════════════════════════════════
\section{In laboratorio: una connessione vista da vicino}
% ══════════════════════════════════════════════════════════════

Una macchina, due finestre. \texttt{nc} dà un estremo TCP senza scrivere un programma: in ascolto da una parte, in
connessione dall'altra.

\begin{lstlisting}[language=bash, title={Due terminali, una connessione}]
$ nc -l 5001                         # server: in ascolto sulla porta 5001
Maradona
sfogliatella

$ ss -tln 'sport = :5001'            # qualcuno ascolta?
$ nc localhost 5001                  # client: connettersi e scrivere
Maradona
sfogliatella
\end{lstlisting}

Ciò che si scrive sul client compare sul server. Nell'altro verso viaggia solo una conferma, un segmento senza alcun
dato. Per catturare il traffico in Wireshark si usa il filtro \texttt{tcp \&\& tcp.port==5001 \&\& ip}:
\begin{itemize}
    \item tre clausole, una per livello della domanda: il protocollo di trasporto, la porta, il protocollo di rete sotto;
    \item quella centrale basterebbe da sola, perché nient'altro sulla macchina usa la 5001: le altre due rendono esplicita
          l'intenzione;
    \item tutto avviene su \texttt{127.0.0.1}: i due estremi sono lo stesso host, e nemmeno un pacchetto raggiunge un cavo.
          L'interfaccia su cui catturare è quindi quella di \textbf{loopback}; puntando lo sniffer su quella reale non si
          vede niente.
\end{itemize}

\begin{lstlisting}[language=bash, title={Che cosa mostra la cattura}]
$ tshark -i lo -f "tcp port 5001"
 1  0.000000 41230 -> 5001  [SYN]       Seq=0  Win=65495 Len=0
 2  0.000018  5001 -> 41230 [SYN, ACK]  Seq=0  Ack=1     Len=0
 3  0.000027 41230 -> 5001  [ACK]       Seq=1  Ack=1     Len=0
 4  6.213004 41230 -> 5001  [PSH, ACK]  Seq=1  Ack=1     Len=9
 5  6.213021  5001 -> 41230 [ACK]       Seq=1  Ack=10    Len=0
 6  9.402771 41230 -> 5001  [PSH, ACK]  Seq=10 Ack=1     Len=13
 7  9.402788  5001 -> 41230 [ACK]       Seq=1  Ack=23    Len=0
 8 19.884113 41230 -> 5001  [FIN, ACK]  Seq=23 Ack=1     Len=0
\end{lstlisting}

\begin{itemize}
    \item Tre segmenti prima che si prema un tasto: l'handshake.
    \item Poi un segmento per ogni riga scritta, di nove e tredici byte (c'è dentro il carattere di a capo), ciascuno con una
          conferma pura che torna indietro: \texttt{Ack=10} è $1 + 9$, \texttt{Ack=23} è $10 + 13$.
    \item L'ultima riga è il Ctrl-C, che chiude e fa partire il FIN.
    \item I numeri di sequenza partono da 0 perché Wireshark e \texttt{tshark} mostrano numeri \textbf{relativi}: sottraggono
          l'ISN casuale per renderli leggibili. Quelli veri sono nel dettaglio del pacchetto.
\end{itemize}

\begin{lstlisting}[language=bash, title={La stessa connessione, con un file}]
$ nc -l 12338 > real2.doc            # server: scrive quanto arriva
$ nc 127.0.0.1 12338 < real.doc      # client: invia un file
^C
\end{lstlisting}

\begin{itemize}
    \item Le due ridirezioni fanno tutto: \texttt{nc} invia il proprio standard input e stampa ciò che riceve. La copia
          avviene su TCP.
    \item Ciò che viaggia è un flusso di byte senza struttura: niente sul filo dice dove finisce il file.
    \item \texttt{nc} non si ferma quando l'input finisce, quindi il trasferimento si chiude a mano, ed è il Ctrl-C a inviare
          il FIN. Qui il FIN è il segnale di fine file: un protocollo che deve inviare due file su una connessione deve dirlo
          da sé (è la delimitazione per chiusura della parte sulla programmazione).
\end{itemize}

\Needspace{20\baselineskip}
% ══════════════════════════════════════════════════════════════
\section{Riepilogo}
% ══════════════════════════════════════════════════════════════

\riepilogo{
\begin{itemize}
    \item Accordarsi su un canale inaffidabile con certezza è impossibile (\textbf{problema dei due eserciti}): il
          three-way handshake è imperfetto ma adeguato.
    \item \textbf{Apertura}: SYN (Seq = client\_isn), SYNACK (Seq = server\_isn, Ack = client\_isn+1), ACK (Ack =
          server\_isn+1, può contenere dati). Numeri iniziali casuali.
    \item Stati di apertura: CLOSED, LISTEN / SYN SENT, SYN RECEIVED, ESTABLISHED.
    \item \textbf{Chiusura}: FIN, ACK, FIN, ACK (ACK e FIN possono essere uniti); stati FIN WAIT 1/2, CLOSE WAIT,
          LAST ACK, CLOSING, TIME WAIT, CLOSED.
    \item \texttt{listen} crea due code (mezze aperte e completate); \texttt{accept} non partecipa all'handshake. I SYN cookie
          difendono dal SYN flood senza tenere stato.
    \item Un FIN chiude una direzione; lo 0 di \texttt{recv} è il FIN del peer. Un CLOSE WAIT che non finisce è un bug
          nostro; TIME WAIT protegge l'ultima conferma e i quattro numeri, ed è il motivo di \texttt{SO\_REUSEADDR}.
    \item I \textbf{timer} risolvono le perdite residue; RST rifiuta o interrompe le connessioni.
    \item Resta aperto: la finestra di ricezione nell'header, e che cosa fa un mittente quando il limite non è il ricevente
          ma la rete.
\end{itemize}
}
